\chapter{Warunkowa wartość oczekiwana i projekcje w \texorpdfstring{$L^2$}{L2}}
\label{ch:condexp}
\paragraph{Odsyłacze międzytomowe.} Statystyczne zastosowania warunkowej wartości oczekiwanej to teoria Rao--Blackwella i reguły Bayesa w tomie IV, rozdziałach~\ref{IV-ch:sufficiency} i~\ref{IV-ch:bayes-foundations}.

\section{Dlaczego warunkowanie wymaga nowej definicji?}

W elementarnych zadaniach, gdy $B\in\mathcal F$ i
$\Prob(B)>0$, używaliśmy
$\E[X\mid B]=\E[X\ind B]/\Prob(B)$.
Ten zapis opisuje jedną liczbę: średnią wartości $X$
po ograniczeniu próby do $B$. W wielu zastosowaniach
obserwujemy jednak nie jedno zdarzenie, lecz pewną
\emph{informację}: wynik części doświadczenia, wartość
pomiaru lub przynależność do jednego z wielu zbiorów.
Zanim poznamy tę informację, średnia zależna od niej
jest zmienną losową, a nie ustaloną liczbą.

Informację reprezentuje pod-$\sigma$-ciało
$\mathcal G\subseteq\mathcal F$: zdarzenia z
$\mathcal G$ są dokładnie tymi, których zajście
potrafimy rozstrzygnąć na podstawie dostępnej obserwacji.
Wartość $\E[X\mid\mathcal G]$ musi zatem
spełniać dwa warunki. Po pierwsze ma zależeć jedynie
od dostępnej informacji, czyli być
$\mathcal G$-mierzalna. Po drugie na \emph{każdym}
zdarzeniu $A\in\mathcal G$ jej całka powinna być
taka sama jak całka oryginalnej zmiennej $X$.
Ten drugi warunek wyznacza właściwą średnią bez
odwoływania się do nieformalnego „warunkowania
na wynik o zerowym prawdopodobieństwie”.

\begin{definicja}[Warunkowa wartość oczekiwana]
Niech $(\Omega,\mathcal F,\Prob)$ będzie przestrzenią
probabilistyczną, $\mathcal G\subseteq\mathcal F$
niech będzie pod-$\sigma$-ciałem, a $X\in L^1(\Omega,
\mathcal F,\Prob)$. Zmienną losową
$Y\in L^1(\Omega,\mathcal G,\Prob)$ nazywamy
\emph{warunkową wartością oczekiwaną $X$ względem
$\mathcal G$}, gdy
\[
\int_A Y\,\dd\Prob=\int_A X\,\dd\Prob
\qquad\text{dla każdego }A\in\mathcal G.
\]
Piszemy $Y=\E[X\mid\mathcal G]$.
Równość rozumiemy z dokładnością do zbioru
$\Prob$-miary zero.
\end{definicja}

\begin{twierdzenie}[Istnienie i jednoznaczność przez Radona--Nikodyma]
Dla każdej przestrzeni probabilistycznej
$(\Omega,\mathcal F,\Prob)$, pod-$\sigma$-ciała
$\mathcal G\subseteq\mathcal F$ i zmiennej
$X\in L^1(\Prob)$ istnieje warunkowa wartość
oczekiwana $\E[X\mid\mathcal G]$, która jest
jednoznaczna $\Prob$-prawie na pewno.
\end{twierdzenie}
\begin{proof}
Rozłóżmy $X=X^+-X^-$, gdzie
$X^+=\max(X,0)$ i $X^-=\max(-X,0)$.
Na $(\Omega,\mathcal G)$ definiujemy skończone
miary dodatnie $\nu_\pm(A)=\E[X^\pm\ind A]$.
Są przeliczalnie addytywne przez twierdzenie
Tonellego i spełniają $\nu_\pm\ll\Prob|_\mathcal G$:
gdy $\Prob(A)=0$, to $\nu_\pm(A)=0$.
Twierdzenie Radona--Nikodyma, zastosowane do miar
na \emph{mniejszym} $\sigma$-ciele $\mathcal G$,
daje nieujemne funkcje $Y_\pm$, mierzalne względem
$\mathcal G$, takie że
$\nu_\pm(A)=\int_A Y_\pm\,\dd\Prob$ dla
każdego $A\in\mathcal G$. Ponieważ
$\E Y_\pm=\nu_\pm(\Omega)=\E X^\pm<\infty$,
zmienna $Y=Y_+-Y_-$ należy do $L^1$ i spełnia
wymaganą tożsamość.

Jeśli $Y$ i $Z$ spełniają tę samą tożsamość,
różnica $D=Y-Z$ jest $\mathcal G$-mierzalna i
$\int_A D\,\dd\Prob=0$ dla wszystkich
$A\in\mathcal G$. W szczególności dla
$A_+=\{D>0\}$ całka $\E[D\ind{A_+}]=0$.
Nieujemna zmienna $D\ind{A_+}$ jest więc zerowa
p.n. Podobnie dla $A_-=\{D<0\}$.
Otrzymujemy $Y=Z$ p.n. Istnienie i jednoznaczność
wynikają zatem dokładnie z konstrukcji pochodnej
Radona--Nikodyma.
\end{proof}

\begin{uwaga}
Definicja nie wymaga, aby wszystkie zdarzenia
z $\mathcal G$ miały dodatnie prawdopodobieństwo.
Wartości warunkowej wartości oczekiwanej na zbiorach
zerowych nie są wyznaczone jednoznacznie. Jest to
własność definicji, a nie przeszkoda techniczna.
\end{uwaga}

\section{Pierwsze modele: od średniej liczbowej do funkcji obserwacji}

\begin{przyklad}[Brak informacji i pełna informacja]
Jeżeli $\mathcal G=\{\varnothing,\Omega\}$, to
każda $\mathcal G$-mierzalna zmienna rzeczywista jest
stała. Warunek całkowy wymusza
$\E[X\mid\mathcal G]=\E X$ p.n. Gdy
$\mathcal G=\mathcal F$, warunek spełnia sama
zmienna $X$, stąd $\E[X\mid\mathcal F]=X$ p.n.
Warunkowanie interpoluje między całkowitą niewiedzą
i pełną znajomością wyniku.
\end{przyklad}

\begin{twierdzenie}[Średnia na atomach skończonego podziału]
Niech $A_1,\ldots,A_m\in\mathcal F$ tworzą
rozłączny podział $\Omega$ i niech
$\mathcal G=\sigma(A_1,\ldots,A_m)$.
Jeśli $X\in L^1$, to wersją $\E[X\mid\mathcal G]$ jest
\[
\sum_{j:\Prob(A_j)>0}
\frac{\E[X\ind{A_j}]}{\Prob(A_j)}\ind{A_j},
\]
przy dowolnych wartościach na atomach zerowych.
\end{twierdzenie}
\begin{proof}
Zmienna z tezy jest stała na każdym atomie $A_j$,
więc jest $\mathcal G$-mierzalna. Na atomie
z $\Prob(A_j)>0$ jej całka jest równa
$\E[X\ind{A_j}]$; na atomie zerowym obie całki
wynoszą zero. Każde zdarzenie z $\mathcal G$
jest sumą pewnych atomów, zatem dodając równości
całek dla atomów otrzymujemy równość dla wszystkich
$A\in\mathcal G$.
\end{proof}

\begin{przyklad}[Jedna informacja o wyniku dwóch kostek]
Niech $D_1,D_2$ będą niezależnymi wynikami dwóch
uczciwych kostek. Obserwujemy tylko, czy $D_1$ jest
parzysta, tj. $\mathcal G=\sigma(\{D_1\in\{2,4,6\}\})$.
Dla $X=D_1+D_2$ mamy
$\E[D_1\mid D_1\text{ parzysta}]=4$,
$\E[D_1\mid D_1\text{ nieparzysta}]=3$,
a $\E D_2=7/2$. Zatem
\[
\E[X\mid\mathcal G]
=\frac{15}{2}\ind{\{D_1\text{ parzysta}\}}
+\frac{13}{2}\ind{\{D_1\text{ nieparzysta}\}}.
\]
To dwuwartościowa zmienna losowa: zanim zobaczymy
parzystość, nie wiemy, którą z dwóch średnich otrzymamy.
\end{przyklad}

\begin{przyklad}[Wiedza o wartości bezwzględnej nie określa znaku]
Niech $U\sim\Unif(-1,1)$ i
$\mathcal G=\sigma(U^2)$. Symetria rozkładu sprawia,
że $\E[U\mid\mathcal G]=0$, mimo że $U$ nie jest
zmienną stałą. Aby to uzasadnić bez odwoływania się
do intuicji, zauważmy, że dla każdego borelowskiego
$B\subseteq[0,1]$ mamy
$\E[U\ind{\{U^2\in B\}}]=0$:
całka po symetrycznym zbiorze z funkcji nieparzystej
jest zerowa. Ponieważ zdarzenia $\{U^2\in B\}$
tworzą $\sigma(U^2)$, zmienna zerowa spełnia
definicję warunkowej wartości oczekiwanej.
Jednocześnie $\E[U^2\mid\sigma(U^2)]=U^2$.
\end{przyklad}

\begin{twierdzenie}[Zmienna warunkowa jako funkcja obserwacji]
Jeżeli $V:\Omega\to\R$ jest zmienną losową,
$\mathcal G=\sigma(V)$ oraz $Y$ jest rzeczywistą
$\mathcal G$-mierzalną zmienną losową, to istnieje
funkcja borelowska $g:\R\to\R$ taka, że
$Y=g(V)$ na całej $\Omega$. W szczególności
istnieje funkcja borelowska $g$ spełniająca
$\E[X\mid\sigma(V)]=g(V)$ prawie na pewno.
\end{twierdzenie}
\begin{proof}
Każde zdarzenie $A\in\sigma(V)$ ma postać
$V^{-1}(B)$ dla pewnego $B\in\mathcal B(\R)$,
ponieważ rodzina wszystkich takich przeciwobrazów
jest dokładnie $\sigma(V)$. Dla prostej funkcji
$Y=\sum_j c_j\ind{A_j}$ otrzymujemy reprezentację
$Y=g(V)$ z borelowską funkcją
$g=\sum_j c_j\ind{B_j}$.
Dla dowolnej mierzalnej rzeczywistej $Y$
konstruujemy ciąg funkcji prostych $Y_m\to Y$
punktowo (osobno aproksymując części dodatnią
i ujemną). Każdą zapisujemy $Y_m=g_m(V)$.
Funkcja $g(t)=\lim_m g_m(t)$ w punktach, w których
granica jest skończona, oraz $g(t)=0$ gdzie indziej,
jest borelowska. Ponieważ $Y_m(\omega)\to Y(\omega)$
dla każdego $\omega$, dostajemy $g(V(\omega))=Y(\omega)$.
\end{proof}

\section{Reguły rachunkowe i twierdzenie o wieży}

\begin{twierdzenie}[Podstawowe własności warunkowej wartości oczekiwanej]
Niech $X,Y\in L^1(\Omega,\mathcal F,\Prob)$,
$\mathcal G\subseteq\mathcal F$ i $a,b\in\R$.
Wszystkie poniższe równości oraz nierówności zachodzą
prawie na pewno:
\begin{enumerate}
\item \emph{Liniowość:}
$\E[aX+bY\mid\mathcal G]=a\E[X\mid\mathcal G]+b\E[Y\mid\mathcal G]$.
\item \emph{Dodatniość i monotoniczność:} jeśli $X\geq0$,
to $\E[X\mid\mathcal G]\geq0$; jeśli $X\leq Y$,
to $\E[X\mid\mathcal G]\leq\E[Y\mid\mathcal G]$.
\item \emph{Zachowanie całki:}
$\E[\E[X\mid\mathcal G]]=\E X$.
\item \emph{Znane wielkości:} jeśli $X$ jest
$\mathcal G$-mierzalna, to
$\E[X\mid\mathcal G]=X$.
\item \emph{Wyciąganie czynnika:} jeśli $H$ jest
ograniczoną $\mathcal G$-mierzalną zmienną, to
$\E[HX\mid\mathcal G]=H\E[X\mid\mathcal G]$.
\item \emph{Niezależność:} jeśli $X$ jest niezależna
od $\mathcal G$ (tj. od każdego zdarzenia
w $\mathcal G$), to $\E[X\mid\mathcal G]=\E X$.
\end{enumerate}
\end{twierdzenie}
\begin{proof}
Liniowość wynika z liniowości całki na każdym
$A\in\mathcal G$ oraz jednoznaczności definicji.
Dla dodatniości niech $Z=\E[X\mid\mathcal G]$
i $X\geq0$. Gdyby $\Prob(Z<0)>0$, to dla
pewnego $r\in\N$ zdarzenie $A=\{Z\leq-1/r\}$
miałoby dodatnią miarę. Wówczas
$\int_A Z\,\dd\Prob\leq-\Prob(A)/r<0$,
co przeczy $\int_AZ\,\dd\Prob=\int_AX\,\dd\Prob\geq0$.
Monotoniczność otrzymujemy przez zastosowanie
dodatniości do $Y-X$.

Własność (3) to definicja dla $A=\Omega$.
Własność (4) wynika z tego, że sama zmienna $X$
spełnia oba warunki definicji. Dla (5) zaczynamy
od $H=\ind B$, $B\in\mathcal G$: dla
$A\in\mathcal G$ mamy
$\int_A\ind B\,\E[X\mid\mathcal G] \,\dd\Prob
=\int_{A\cap B}X\,\dd\Prob$.
Liniowość rozszerza to na funkcje proste, a
jednostajna aproksymacja ograniczonych funkcji
mierzalnych przez proste oraz zdominowana zbieżność
na dowolne ograniczone $H$.
Wreszcie dla niezależnego $X$ i $A\in\mathcal G$
$\E[X\ind A]=\E X\,\Prob(A)$.
Dla całkowalnych zmiennych wynika to z definicji
niezależności i aproksymacji funkcjami prostymi;
stała $\E X$ spełnia zatem wymaganą tożsamość.
\end{proof}

\begin{twierdzenie}[Własność wieży]
Jeśli $\mathcal G\subseteq\mathcal H\subseteq\mathcal F$
są $\sigma$-ciałami i $X\in L^1$, to
\[
\E[\E[X\mid\mathcal H]\mid\mathcal G]
=\E[X\mid\mathcal G],
\qquad
\E[\E[X\mid\mathcal G]\mid\mathcal H]
=\E[X\mid\mathcal G].
\]
\end{twierdzenie}
\begin{proof}
Pierwsza lewa strona jest $\mathcal G$-mierzalna.
Dla $A\in\mathcal G\subseteq\mathcal H$
na mocy dwóch kolejnych definicji mamy
\[
\int_A \E[\E[X\mid\mathcal H]\mid\mathcal G]
\,\dd\Prob
=\int_A\E[X\mid\mathcal H]\,\dd\Prob
=\int_A X\,\dd\Prob.
\]
Jednoznaczność daje pierwszą równość.
Dla drugiej zauważamy, że $\E[X\mid\mathcal G]$
jest już $\mathcal H$-mierzalna, więc jej ponowne
warunkowanie względem $\mathcal H$ nic nie zmienia.
\end{proof}

\begin{przyklad}[Trzy poziomy informacji]
Niech $D_1,D_2$ będą niezależnymi wynikami kostek,
$X=D_1+D_2$, $\mathcal H=\sigma(D_1)$,
a $\mathcal G=\sigma(\ind{\{D_1\text{ parzysta}\}})$.
Wtedy
$\E[X\mid\mathcal H]=D_1+7/2$,
a następnie $\E[D_1+7/2\mid\mathcal G]$
daje tę samą dwuwartościową zmienną co
bezpośrednie $\E[X\mid\mathcal G]$ z wcześniejszego
przykładu. Własność wieży oznacza, że można
najpierw uśrednić względem informacji dokładniejszej,
a potem grubszej; wynik nie zależy od tego
pośredniego kroku.
\end{przyklad}

\begin{przyklad}[Wieża nie działa dowolnie dla nieporównywalnych informacji]
Niech $\Omega=\{1,2,3\}$, każdy punkt z masą $1/3$,
$X=\ind{\{1\}}$, $\mathcal G=\sigma(\{1\})$
i $\mathcal H=\sigma(\{2\})$.
Zapisując zmienne jako wektory wartości na
$(1,2,3)$, dostajemy
\[
\E[\E[X\mid\mathcal G]\mid\mathcal H]
=(1/2,0,1/2),\qquad
\E[\E[X\mid\mathcal H]\mid\mathcal G]
=(1/2,1/4,1/4).
\]
Te wyniki są różne. Własność wieży zakłada
\emph{zawieranie się} odpowiednich $\sigma$-ciał;
bez niego kolejne projekcje na różne zbiory
informacji nie muszą być przemienne.
\end{przyklad}

\section{Granice pod warunkowaniem i nierówność Jensena}

\begin{twierdzenie}[Monotoniczna i zdominowana zbieżność warunkowa]
Niech $\mathcal G\subseteq\mathcal F$.
\begin{enumerate}
\item Jeśli $0\leq X_n\uparrow X$ p.n., a $X\in L^1$,
to $\E[X_n\mid\mathcal G]\uparrow\E[X\mid\mathcal G]$
prawie na pewno.
\item Jeśli $X_n\to X$ p.n., a $|X_n|\leq W$
p.n. dla wszystkich $n$, gdzie $W\in L^1$, to
$\E[X_n\mid\mathcal G]\to\E[X\mid\mathcal G]$
prawie na pewno oraz w $L^1$.
\end{enumerate}
\end{twierdzenie}
\begin{proof}
W (1) monotoniczność warunkowej wartości oczekiwanej
daje niemalejący ciąg $Y_n=\E[X_n\mid\mathcal G]$.
Jest on ograniczony z góry przez
$\E[X\mid\mathcal G]\in L^1$, więc granica
$Y=\lim_nY_n$ jest skończona p.n. Dla dowolnego
$A\in\mathcal G$ twierdzenie o zbieżności monotonicznej
daje
$\E[Y\ind A]=\lim_n\E[Y_n\ind A]
=\lim_n\E[X_n\ind A]=\E[X\ind A]$.
Jednoznaczność dowodzi (1).

Dla (2) mamy $X_n+W\geq0$ i $W-X_n\geq0$,
więc można zastosować (1) po uprzednim zastąpieniu
ogólnego ciągu przez suprema i infima ogonowe.
Dokładniej, niech
$L_n=\inf_{k\geq n}X_k$ i
$U_n=\sup_{k\geq n}X_k$. Wtedy
$L_n\uparrow X$, $U_n\downarrow X$ p.n., a
$-W\leq L_n\leq X_n\leq U_n\leq W$.
Wersja monotoniczna zastosowana do $L_n+W$ oraz
do $W-U_n$ daje
$\E[L_n\mid\mathcal G]\uparrow\E[X\mid\mathcal G]$
i $\E[U_n\mid\mathcal G]\downarrow\E[X\mid\mathcal G]$.
Monotoniczność i twierdzenie o trzech ciągach
dają zbieżność p.n. warunkowych wartości $X_n$.
Są one ograniczone modułem przez $\E[W\mid\mathcal G]$
na mocy dodatniości, a ten majorant jest całkowalny;
zdominowana zbieżność daje też zbieżność w $L^1$.
\end{proof}

\begin{twierdzenie}[Nierówność Jensena warunkowa]
Niech $X\in L^1(\Prob)$, a $\phi:\R\to\R$ będzie
skończoną funkcją wypukłą taką, że
$\phi(X)\in L^1(\Prob)$. Wówczas dla każdego
$\mathcal G\subseteq\mathcal F$
zachodzi
\[
\phi(\E[X\mid\mathcal G])
\leq\E[\phi(X)\mid\mathcal G]
\qquad\text{prawie na pewno}.
\]
\end{twierdzenie}
\begin{proof}
Dla dowolnego $r\in\mathbb Q$ wybierzmy prostą
podpierającą wykres wypukłej funkcji $\phi$ w punkcie
$r$, czyli liczbę $s_r$ taką, że
$\phi(x)\geq\phi(r)+s_r(x-r)$ dla wszystkich $x\in\R$.
Istnienie $s_r$ wynika z nierówności pomiędzy
lewą i prawą pochodną funkcji wypukłej.
Dodatniość i liniowość warunkowej wartości
oczekiwanej dają
\[
\E[\phi(X)\mid\mathcal G]
\geq\phi(r)+s_r(\E[X\mid\mathcal G]-r)
\quad\text{p.n.}
\]
Ponieważ liczb wymiernych jest przeliczalnie wiele,
możemy wybrać \emph{jeden} zbiór pełnej miary,
na którym nierówność zachodzi dla wszystkich $r$.
Dla ciągłej skończonej funkcji wypukłej zachodzi
$\phi(x)=\sup_{r\in\mathbb Q}\{\phi(r)+s_r(x-r)\}$:
z prawej strony mamy co najwyżej $\phi(x)$,
a dla wymiernych $r\to x$ wyrażenia zbliżają się
do $\phi(x)$, ponieważ lokalnie nachylenia
podpierające są ograniczone. Podstawienie
$x=\E[X\mid\mathcal G]$ daje tezę.
\end{proof}

\begin{wniosek}[Kontrakcja w $L^p$]
Niech $1\leq p<\infty$, $X\in L^p(\Prob)$ i
$\mathcal G\subseteq\mathcal F$. Wówczas
$\E[X\mid\mathcal G]\in L^p$ oraz
$\|\E[X\mid\mathcal G]\|_p\leq\|X\|_p$.
\end{wniosek}
\begin{proof}
Dla $p\geq1$ funkcja $x\mapsto|x|^p$ jest wypukła.
Stosując warunkowego Jensena, a następnie całkując,
dostajemy
$\E|\E[X\mid\mathcal G]|^p\leq
\E\E[|X|^p\mid\mathcal G]=\E|X|^p$.
\end{proof}

\begin{przyklad}[Dlaczego średnia warunkowa zmniejsza błąd średniokwadratowy]
Niech $X=Z+\varepsilon$, gdzie $Z$ jest całkowalna
i $\varepsilon$ jest niezależna od $Z$ o średniej
zero i skończonej wariancji. Jeśli obserwujemy $Z$,
to $\E[X\mid\sigma(Z)]=Z$. Zmienna $Z$ jest więc
średnią przewidywaną na podstawie obserwacji, a
$\varepsilon$ pozostaje nieprzewidywalnym błędem.
Gdy oba składniki są w $L^2$, ich zależność
z formalną projekcją zostanie wyprowadzona poniżej.
\end{przyklad}

\section{Interpretacja geometryczna: rzut ortogonalny w \texorpdfstring{$L^2$}{L2}}

Oznaczmy przez $L^2(\mathcal F)$ przestrzeń
rzeczywistych zmiennych losowych o skończonym drugim
momencie, utożsamiając zmienne równe prawie na pewno.
Iloczyn skalarny $\langle U,V\rangle=\E[UV]$
czyni tę przestrzeń przestrzenią Hilberta.
Podprzestrzeń $L^2(\mathcal G)$ składa się z klas
$\mathcal G$-mierzalnych zmiennych całkowalnych z
kwadratem. Jest domknięta w $L^2(\mathcal F)$:
jeżeli $Y_n\in L^2(\mathcal G)$ i $Y_n\to Y$
w $L^2$, to z kryterium podciągowego istnieje
podciąg zbiegający do $Y$ prawie na pewno.
Jego granica jest $\mathcal G$-mierzalna (po wyborze
odpowiedniej wersji $Y$), co dowodzi domkniętości.

\begin{twierdzenie}[Warunkowa wartość oczekiwana jako rzut]
Niech $X\in L^2(\Omega,\mathcal F,\Prob)$ i
$\mathcal G\subseteq\mathcal F$. Wtedy
$P_{\mathcal G}X:=\E[X\mid\mathcal G]$ jest
jedynym elementem $L^2(\mathcal G)$ spełniającym
\[
\E\big[(X-P_{\mathcal G}X)Z\big]=0
\qquad\text{dla każdego } Z\in L^2(\mathcal G).
\]
W szczególności dla każdego $Z\in L^2(\mathcal G)$
zachodzi tożsamość Pitagorasa
\[
\E|X-Z|^2
=\E|X-P_{\mathcal G}X|^2
+\E|P_{\mathcal G}X-Z|^2.
\]
Zmienna $P_{\mathcal G}X$ jest zatem jedynym
najlepszym przybliżeniem $X$ przez zmienne
$\mathcal G$-mierzalne w normie $L^2$.
\end{twierdzenie}
\begin{proof}
Kontrakcja w $L^2$ dowodzi, że $P_{\mathcal G}X$
należy do $L^2(\mathcal G)$. Z definicji mamy
$\E[(X-P_{\mathcal G}X)\ind A]=0$
dla wszystkich $A\in\mathcal G$.
Liniowość rozszerza równość na proste funkcje
$\mathcal G$-mierzalne. Jeżeli $Z\in L^2(\mathcal G)$,
istnieją proste $Z_j\to Z$ w $L^2$. Na mocy
nierówności Cauchy'ego--Schwarza
\[
\left|\E[(X-P_{\mathcal G}X)(Z_j-Z)]\right|
\leq\|X-P_{\mathcal G}X\|_2\,\|Z_j-Z\|_2\to0.
\]
Przechodzimy do granicy i otrzymujemy ortogonalność.

Teraz rozkładamy
$X-Z=(X-P_{\mathcal G}X)+(P_{\mathcal G}X-Z)$.
Drugi składnik należy do $L^2(\mathcal G)$,
więc wyraz mieszany w kwadracie normy jest zerowy.
To daje tożsamość Pitagorasa. Drugi składnik
po prawej jest nieujemny i znika wtedy i tylko
wtedy, gdy $Z=P_{\mathcal G}X$ w $L^2$.
Jednoznaczność najlepszego przybliżenia jest dowiedziona.
\end{proof}

\begin{wniosek}[Własności operatora projekcji]
Dla $X,Y\in L^2(\mathcal F)$ i
$\mathcal G\subseteq\mathcal F$ operator
$P_{\mathcal G}:L^2(\mathcal F)\to L^2(\mathcal F)$
jest liniowy, samosprzężony, idempotentny
($P_{\mathcal G}^2=P_{\mathcal G}$) i kontrakcyjny.
Jeśli $\mathcal G\subseteq\mathcal H$, to
$P_{\mathcal G}P_{\mathcal H}
=P_{\mathcal H}P_{\mathcal G}=P_{\mathcal G}$.
\end{wniosek}
\begin{proof}
Liniowość i kontrakcję wykazano wcześniej,
a idempotentność wynika z warunkowania zmiennej
już $\mathcal G$-mierzalnej. Z ortogonalności
$X-P_{\mathcal G}X$ do $L^2(\mathcal G)$
dostajemy, dla dowolnych $X,Y$,
$\langle P_{\mathcal G}X,Y\rangle
=\langle P_{\mathcal G}X,P_{\mathcal G}Y\rangle
=\langle X,P_{\mathcal G}Y\rangle$,
co jest samosprzężonością.
Ostatnie równości to dokładnie obie wersje
własności wieży dla zagnieżdżonych $\sigma$-ciał.
\end{proof}

\begin{przyklad}[Rzut na stałe i elementarna tożsamość wariancyjna]
Jeśli $\mathcal G=\{\varnothing,\Omega\}$ i $X\in L^2$,
to $P_{\mathcal G}X=\E X$. Równość Pitagorasa
przy $Z=c\in\R$ daje
$\E(X-c)^2=\Var(X)+(\E X-c)^2$.
Wyjaśnia to geometrycznie, dlaczego \emph{średnia}
minimalizuje oczekiwany kwadrat błędu, a nie np.
mediana (ta ostatnia minimalizuje inny rodzaj straty).
\end{przyklad}

\section{Wariancja warunkowa, dekompozycja błędu i kowariancja}

\begin{definicja}
Dla $X\in L^2(\Prob)$ i pod-$\sigma$-ciała
$\mathcal G\subseteq\mathcal F$ definiujemy
\emph{warunkową wariancję}
\[
\Var(X\mid\mathcal G)
:=\E[(X-\E[X\mid\mathcal G])^2\mid\mathcal G].
\]
Jest to nieujemna, $\mathcal G$-mierzalna
zmienna losowa.
\end{definicja}

\begin{twierdzenie}[Wzór na wariancję całkowitą]
Dla każdego $X\in L^2(\Prob)$ oraz
$\mathcal G\subseteq\mathcal F$ zachodzą
\[
\Var(X\mid\mathcal G)=
\E[X^2\mid\mathcal G]-\E[X\mid\mathcal G]^2
\quad\text{p.n.}
\]
oraz
\[
\Var(X)=\E[\Var(X\mid\mathcal G)]
+\Var(\E[X\mid\mathcal G]).
\]
\end{twierdzenie}
\begin{proof}
Oznaczmy $M=\E[X\mid\mathcal G]$.
Rozwinięcie kwadratu i reguła wyciągania
czynnika (najpierw dla ograniczonych obcięć $M$,
a potem przez zbieżność w $L^1$ wynikającą
z Cauchy'ego--Schwarza) dają
$\E[X M\mid\mathcal G]=M\E[X\mid\mathcal G]=M^2$.
Podobnie $\E[M^2\mid\mathcal G]=M^2$.
Stąd pierwszy wzór.

Zastosujmy następnie Pitagorasa do rozkładu
$X-\E X=(X-M)+(M-\E X)$.
Jego pierwszy składnik jest ortogonalny do
drugiego, bo drugi jest $\mathcal G$-mierzalny.
Otrzymujemy
$\Var(X)=\E(X-M)^2+\E(M-\E X)^2$.
Z własności zachowania wartości oczekiwanej
$\E(M)=\E X$, a z definicji wariancji
warunkowej $\E(X-M)^2=\E\Var(X\mid\mathcal G)$.
To dowodzi drugiej równości.
\end{proof}

\begin{przyklad}[Losowanie z dwóch grup]
Wybieramy grupę $G\in\{0,1\}$ z równymi
prawdopodobieństwami, a następnie obserwujemy
$X=2G+Z$, gdzie $Z$ jest niezależna od $G$,
$\E Z=0$, $\Var(Z)=1$. Wówczas
$\E[X\mid G]=2G$ i
$\Var(X\mid G)=1$. Wzór na wariancję całkowitą
daje $\Var(X)=1+\Var(2G)=1+1=2$.
Pierwszy składnik opisuje rozproszenie \emph{wewnątrz}
grup, drugi różnicę pomiędzy średnimi grupowymi.
\end{przyklad}

\begin{twierdzenie}[Kowariancja całkowita]
Niech $X,Y\in L^2(\Prob)$ i
$\mathcal G\subseteq\mathcal F$. Wówczas
\[
\Cov(X,Y)=\E[\Cov(X,Y\mid\mathcal G)]
+\Cov(\E[X\mid\mathcal G],\E[Y\mid\mathcal G]),
\]
gdzie
$\Cov(X,Y\mid\mathcal G)=
\E[XY\mid\mathcal G]-
\E[X\mid\mathcal G]\E[Y\mid\mathcal G]$.
\end{twierdzenie}
\begin{proof}
Niech $M=\E[X\mid\mathcal G]$ i
$N=\E[Y\mid\mathcal G]$. Własność wieży daje
$\E\Cov(X,Y\mid\mathcal G)=\E[XY]-\E[MN]$.
Ponieważ $\E M=\E X$ i $\E N=\E Y$,
$\Cov(M,N)=\E[MN]-\E X\E Y$.
Dodajemy obie równości i otrzymujemy tezę.
\end{proof}

\section{Warunkowanie względem zmiennej ciągłej i gęstości}

Jeżeli znamy wspólną gęstość pary $(X,V)$,
możemy jawnie obliczyć funkcję $g$ z reprezentacji
$\E[X\mid\sigma(V)]=g(V)$. Należy przy tym
oddzielić równość całkową określającą średnią
od naiwnego warunku „$V=v$”, który w modelu
ciągłym zwykle ma prawdopodobieństwo zero.

\begin{twierdzenie}[Wzór przez gęstość warunkową]
Niech $(X,V)$ będzie rzeczywistym wektorem losowym
o wspólnej gęstości $f_{X,V}$ względem miary
$\lambda_2$ oraz niech $\E|X|<\infty$.
Niech $f_V(v)=\int_\R f_{X,V}(x,v)\,\dd\lambda_1(x)$.
Dla tych $v$, dla których $0<f_V(v)<\infty$ oraz
$\int_\R |x|f_{X,V}(x,v)\,\dd\lambda_1(x)<\infty$,
połóżmy
\[
g(v)=\frac{\int_\R x f_{X,V}(x,v)\,\dd\lambda_1(x)}{f_V(v)},
\]
a w pozostałych punktach niech $g(v)=0$.
Wówczas $g(V)$ jest wersją
$\E[X\mid\sigma(V)]$.
\end{twierdzenie}
\begin{proof}
Z twierdzenia Tonellego i założenia $\E|X|<\infty$
wynika
\[
\int_\R\int_\R |x|f_{X,V}(x,v)
\,\dd\lambda_1(x)\,\dd\lambda_1(v)<\infty.
\]
Stąd licznik w definicji $g$ istnieje dla
$\lambda_1$-prawie wszystkich $v$; zbiory, na których
$f_V=0$ lub $f_V=\infty$, mają zerową masę
względem rozkładu $V$. Ponadto
$\int |g(v)|f_V(v)\,\dd\lambda_1(v)\leq\E|X|$,
więc $g(V)\in L^1$. Dla każdego zbioru borelowskiego
$B\subseteq\R$ twierdzenie Fubiniego daje
\[
\E[g(V)\ind{\{V\in B\}}]
=\int_Bg(v)f_V(v)\,\dd\lambda_1(v)
=\int_B\int_\R xf_{X,V}(x,v)
\,\dd\lambda_1(x)\,\dd\lambda_1(v)
=\E[X\ind{\{V\in B\}}].
\]
Każde zdarzenie z $\sigma(V)$ ma postać
$\{V\in B\}$, więc jest to dokładnie
tożsamość definiująca warunkową wartość oczekiwaną.
\end{proof}

\begin{przyklad}[Jednostajny punkt trójkąta]
Niech $(X,V)$ będzie jednostajny na zbiorze
$T=\{(x,v):0<v<x<1\}$, a więc jego gęstość
wynosi $2\ind T$. Dla $v\in(0,1)$ mamy
$f_V(v)=\int_v^1 2\,\dd\lambda_1(x)=2(1-v)$,
a gęstość $X$ warunkowo przy $V=v$ jest
jednostajna na $(v,1)$. Otrzymujemy
$\E[X\mid V=v]=(1+v)/2$ dla $v\in(0,1)$
i w konsekwencji
$\E[X\mid\sigma(V)]=(1+V)/2$ p.n.
Sprawdźmy to bezpośrednio: dla dowolnego
borelowskiego $B\subseteq(0,1)$
\[
\E[X\ind{\{V\in B\}}]
=2\int_B\!\int_v^1 x\,\dd\lambda_1(x)
\,\dd\lambda_1(v)
=\int_B(1-v^2)\,\dd\lambda_1(v),
\]
podczas gdy
$\E[((1+V)/2)\ind{\{V\in B\}}]
=\int_B (1+v)(1-v)\,\dd\lambda_1(v)$.
Obie całki są równe. Ten rachunek pokazuje
zarówno funkcję warunkową, jak i poprawne znaczenie
warunkowania na konkretną wartość zmiennej ciągłej.
\end{przyklad}

\begin{przyklad}[Suma niezależnych, jednakowo rozłożonych składników]
Niech $X,Y\in L^1$ będą niezależne i jednakowo
rozłożone, $S=X+Y$. Wtedy $\sigma(S)$
nie rozróżnia zamiany $X$ z $Y$. Formalnie
wspólny rozkład $(X,Y)$ jest niezmienniczy na taką
zamianę, więc dla każdego ograniczonego
borelowskiego $h$ zachodzi
$\E[Xh(S)]=\E[Yh(S)]$.
Dodając równości, dostajemy
$2\E[Xh(S)]=\E[Sh(S)]$. Definicja warunkowa
daje zatem $\E[X\mid\sigma(S)]=S/2$,
niezależnie od tego, czy rozkład składników ma gęstość.
\end{przyklad}

\begin{uwaga}[Wersje i zdarzenia zerowe]
Funkcja $g(v)$ w twierdzeniu o gęstości warunkowej
jest określona jednoznacznie tylko dla
$\Law(V)$-prawie wszystkich $v$. Nie ma żadnego
aksjomatu nakazującego konkretną wartość w punkcie,
w którym $f_V(v)=0$. Sformalizowane rozkłady
warunkowe jako rodziny miar zależnych od $v$
można konstruować na standardowych przestrzeniach
borelowskich, ale ta dodatkowa teoria nie jest
potrzebna do definicji i własności
$\E[X\mid\sigma(V)]$.
\end{uwaga}

\section{Co warto zapamiętać przed procesami stochastycznymi?}
Warunkowa wartość oczekiwana działa w dwóch językach.
W języku miary zachowuje całki na zdarzeniach
obserwowalnych. W języku geometrii jest rzutem
ortogonalnym na domkniętą przestrzeń zmiennych
zależnych od dostępnej informacji. Zmiana informacji
zmienia rzut, a własność wieży kontroluje kolejne
etapy ujawniania danych. W drugim tomie właśnie
ten mechanizm będzie podstawą definicji martyngału,
w którym zachodzi $\E[M_{n+1}\mid\mathcal F_n]=M_n$.

\section{Warunkowanie: czynnik nie musi być ograniczony}

W rozdziale~\ref{ch:condexp} reguła wyciągania czynnika została wykazana dla
ograniczonej zmiennej $H$. To założenie nie jest konieczne, ale jego
usunięcie wymaga całkowalności odpowiednich iloczynów.

\begin{twierdzenie}[Wyciąganie czynnika w pełnej wersji]
Niech $(\Omega,\mathcal F,\Prob)$ będzie przestrzenią probabilistyczną,
$\mathcal G\subseteq\mathcal F$, $X\in L^1(\Prob)$ i niech $H$
będzie rzeczywistą $\mathcal G$-mierzalną zmienną losową. Jeżeli
$HX\in L^1(\Prob)$, to $H\E[X\mid\mathcal G]\in L^1(\Prob)$ oraz
\[
\E[HX\mid\mathcal G]=H\E[X\mid\mathcal G]\quad\text{p.n.}
\]
\end{twierdzenie}
\begin{proof}
Z warunkowego Jensena mamy
$|\E[X\mid\mathcal G]|\leq \E[|X|\mid\mathcal G]$.
Dla $H_m=\min(|H|,m)$, który jest nieujemny, ograniczony
i $\mathcal G$-mierzalny, reguła wyciągania czynnika daje
$\E[H_m|X|\mid\mathcal G]=H_m\E[|X|\mid\mathcal G]$.
Po scałkowaniu i przejściu do granicy przez twierdzenie
Beppo Leviego otrzymujemy
\[
\E[|H|\E(|X|\mid\mathcal G)]=\E|HX|<\infty.
\]
W szczególności $H\E[X\mid\mathcal G]$ jest całkowalna.
Dla $H^{(m)}=(-m)\vee(H\wedge m)$ zachodzi
$H^{(m)}X\to HX$ i $H^{(m)}\E[X\mid\mathcal G]\to
H\E[X\mid\mathcal G]$ prawie na pewno, a obie rodziny
są zdominowane odpowiednio przez $|HX|$ oraz
$|H|\E[|X|\mid\mathcal G]$. Stosujemy wcześniejszą
regułę dla $H^{(m)}$, całkujemy po dowolnym $A\in\mathcal G$
i przechodzimy do granicy. Jednoznaczność warunkowej
wartości oczekiwanej kończy dowód.
\end{proof}

\begin{przyklad}[Dlaczego potrzebujemy warunku całkowalności]
Jeśli $X\equiv1$ oraz $H\geq0$ jest $\mathcal G$-mierzalna,
ale $\E H=\infty$, to formalna równość
$\E[HX\mid\mathcal G]=H\E[X\mid\mathcal G]$
nie jest równością dwóch elementów $L^1$.
Można osobno rozwijać warunkowe wartości oczekiwane
nieujemnych zmiennych o nieskończonej średniej,
lecz wymaga to definicji o wartościach w $[0,\infty]$.
\end{przyklad}

\section{Granice i jednostajna całkowalność}

\begin{twierdzenie}[Zbieżność $L^1$ zachowuje się przy warunkowaniu]
Niech $X_n,X\in L^1(\Prob)$, $\mathcal G\subseteq\mathcal F$.
Jeżeli $\E|X_n-X|\to0$, to
\[
\E\bigl|\E[X_n\mid\mathcal G]-\E[X\mid\mathcal G]\bigr|\to0.
\]
W szczególności wyniku tego nie trzeba wyprowadzać
z prawie pewnej zbieżności $X_n$: jest on bezpośrednią
konsekwencją kontrakcji w $L^1$.
\end{twierdzenie}
\begin{proof}
Z liniowości i warunkowej nierówności Jensena
\[
\E\left|\E[X_n-X\mid\mathcal G]\right|
\leq \E\E[|X_n-X|\mid\mathcal G]=\E|X_n-X|\to0.
\]
\end{proof}

\begin{przyklad}[Zbieżność prawie na pewno nie wystarcza]
Na $((0,1),\mathcal B((0,1)),\lambda_1)$ niech
$X_n=n\ind{(0,1/n)}$. Wtedy $X_n\to0$ prawie na pewno,
ale $\E X_n=1$ dla każdego $n$.
Przy trywialnym $\mathcal G$ mamy zatem
$\E[X_n\mid\mathcal G]=1\not\to0=\E[0\mid\mathcal G]$.
To wyjaśnia, dlaczego w twierdzeniu o zbieżności zdominowanej
w rozdziale~\ref{ch:condexp} rzeczywiście potrzebny jest majorant
całkowalny lub inny warunek gwarantujący $L^1$.
\end{przyklad}



\begin{przyklad}[Niezależność warunkowa nie jest niezależnością]
Niech $B_1,B_2$ będą niezależnymi próbami Bernoulliego
z parametrem $1/2$ i $S=B_1+B_2$.
Zmienne $B_1,B_2$ są niezależne, ale po poznaniu
$S=1$ wiemy, że dokładnie jedna z nich jest równa $1$.
Istotnie,
$\Prob(B_1=1,B_2=1\mid S=1)=0$, natomiast
$\Prob(B_1=1\mid S=1)=\Prob(B_2=1\mid S=1)=1/2$.
Iloczyn prawdopodobieństw brzegowych warunkowych
wynosi zatem $1/4$, nie zero. Obserwacja sumy
wprowadza zależność warunkową mimo wcześniejszej
niezależności; w szczególności nie wolno automatycznie
zastępować warunkowej wartości oczekiwanej iloczynu
iloczynem warunkowych wartości oczekiwanych.
\end{przyklad}

\section{Dwa dalsze przykłady warunkowania: rozkład dyskretny i gaussowski}

\begin{przyklad}[Suma niezależnych zmiennych Poissona]
Niech $X\sim\Poi(\lambda)$ i $Y\sim\Poi(\mu)$ będą
niezależne, gdzie $\lambda,\mu>0$, i niech $S=X+Y$.
Z rozkładu splotu wiemy, że $S\sim\Poi(\lambda+\mu)$.
Dla $n\in\N_0$ i $0\leq k\leq n$ obliczamy
\begin{align*}
\Prob(X=k\mid S=n)
&=\frac{\Prob(X=k,Y=n-k)}{\Prob(S=n)}\\
&=\binom nk\left(\frac{\lambda}{\lambda+\mu}\right)^k
\left(\frac{\mu}{\lambda+\mu}\right)^{n-k}.
\end{align*}
Zatem $X$ przy warunku $S=n$ ma rozkład
$\Bin(n,\lambda/(\lambda+\mu))$. Ponieważ
$\sigma(S)$ jest generowana przez rozłączny,
przeliczalny podział $\{S=n\}$,
\[
\E[X\mid\sigma(S)]
=\frac{\lambda}{\lambda+\mu}S,
\quad
\Var(X\mid\sigma(S))
=\frac{\lambda\mu}{(\lambda+\mu)^2}S.
\]
Intuicyjnie suma mówi, ile zdarzeń zaszło,
a parametry $\lambda$ i $\mu$ określają,
jaka część z nich pochodzi z każdego źródła.
\end{przyklad}

\begin{przyklad}[Dwie niezależne obserwacje normalne]
Niech $X,Z\sim\Normal(0,1)$ będą niezależne.
Obserwujemy tylko sumę $S=X+Z$ i chcemy
przewidzieć $X$. Z symetrii otrzymujemy
$\E[X\mid\sigma(S)]=S/2$.
Aby poznać również wariancję warunkową,
połóżmy $D=(X-Z)/2$. Wtedy
$X=S/2+D$, a odwrotną transformacją jest
$X=S/2+D$, $Z=S/2-D$, o wyznaczniku
Jacobianu o module $1$. Wspólna gęstość
$(S,D)$ wynosi
\[
f_{S,D}(s,d)=\frac1{2\pi}
\exp\left(-\frac{s^2}{4}-d^2\right),
\]
co rozkłada się na iloczyn gęstości
$\Normal(0,2)$ dla $S$ i $\Normal(0,1/2)$
dla $D$. Są więc niezależne.
Wobec $X=S/2+D$ i $\E D=0$ mamy
$\Var(X\mid\sigma(S))=\Var(D)=1/2$.
Po obserwacji sumy pozostaje połowa pierwotnej
wariancji pojedynczego składnika.
\end{przyklad}

